\documentclass[10pt,a4paper]{article}

% ---------- Geometry & typography -------------------------------------------
\usepackage[margin=2.2cm]{geometry}
\usepackage[T1]{fontenc}
\usepackage[utf8]{inputenc}
\usepackage{lmodern}
\usepackage{microtype}
\usepackage{csquotes}

% Allow looser line breaks before flagging a paragraph as overfull. Removes
% the small (<40pt) overfull-hbox warnings that arise from long inline
% \texttt{} runs (e.g. "NullProvider/GeminiProvider", parameter lists).
\setlength{\emergencystretch}{3em}
\sloppy

% ---------- Maths -----------------------------------------------------------
\usepackage{amsmath}
\usepackage{amssymb}

% ---------- Tables & figures ------------------------------------------------
\usepackage{booktabs}
\usepackage{array}
\usepackage{graphicx}
\usepackage{xcolor}

% ---------- Code listings ---------------------------------------------------
\usepackage{listings}
\lstset{
  basicstyle=\ttfamily\small,
  breaklines=true,
  columns=fullflexible,
  keywordstyle=\color{blue!60!black},
  commentstyle=\color{black!50},
  stringstyle=\color{green!50!black},
  showstringspaces=false,
  frame=single,
  rulecolor=\color{black!20},
  framerule=0.4pt,
  xleftmargin=4pt,
  xrightmargin=4pt
}

% ---------- Custom macros ---------------------------------------------------
% \keystroke{...} renders a keyboard shortcut (used in the project diary).
\providecommand{\keystroke}[1]{\fbox{\footnotesize\texttt{#1}}}

% ---------- TODO notes for team-filled placeholders -------------------------
\usepackage[colorinlistoftodos,prependcaption,textsize=footnotesize]{todonotes}

% ---------- Hyperlinks (load near the end) ----------------------------------
\usepackage[hidelinks,colorlinks=true,urlcolor=blue!60!black,linkcolor=black,citecolor=blue!50!black]{hyperref}

% ---------- Bibliography ----------------------------------------------------
\usepackage[backend=biber,style=numeric,sorting=nyt,maxbibnames=4]{biblatex}
\addbibresource{references.bib}

% ---------- Document metadata -----------------------------------------------
\title{%
  \textbf{QuantEdge} \\[4pt]
  \large A Backtesting Framework for Retail Quantitative Traders
}
\author{%
  Aaron Arauz \texttt{[Baumender11]} \and
  Luca Di Pietro \texttt{[L-Di-Pietro]} \and
  Stefano Angelo Galdini \texttt{[SAGaldini]} \and
  Filippo Selmi \texttt{[FilippoSelmi]}
}
\date{2026-06-12}

% ============================================================================
\begin{document}
% ============================================================================

\maketitle

\begin{center}
\small
\textbf{USI Lugano \textemdash{} Programming in Finance II (2026)} \\
Project 2.8 \textemdash{} Instructor: P.~H.~Gruber \\
Source repository: \url{https://github.com/L-Di-Pietro/Programming_II_Project_USI} \\
Web app: \url{https://quantedgeusi.up.railway.app/}
\end{center}

\vspace{1em}

\begin{abstract}
\noindent
QuantEdge is a web-based event-driven backtesting framework for retail quantitative
traders, built as the final project for USI Lugano's \emph{Programming in Finance~II}
course. The system targets the two failure modes that cause most retail-backtest losses:
look-ahead bias (architecturally prevented by enforcing bar-$t \to$ bar-$t{+}1$ fills in
the engine) and unrealistic execution assumptions (commission and slippage exposed as
first-class user parameters). It ships eleven trading strategies across five families,
multi-asset support (equities, ETFs, FX, crypto) with native trading calendars per asset
class, both daily and hourly bar resolution, and a six-agent backend in which only two
of the six agents are LLM-backed and are opt-in by default. This document covers the
project plan, development diary, methodology, sample results with an economic
reflection, lessons learned, and an acknowledgement of the generative-AI tools used.
\end{abstract}

\tableofcontents

\vspace{1em}
\hrule
\vspace{1em}

% ---------- Sections --------------------------------------------------------
\section{Project plan}
\label{sec:plan}

\subsection{Problem statement}

Retail quantitative traders lose money for two related reasons. First, many deploy
strategies that have never been rigorously back-tested on historical data. Second,
when they \emph{do} test, they cut corners that systematically inflate apparent
performance: they let the strategy peek at information that would not have been
available in real time (\emph{look-ahead bias}), they neglect transaction costs and
execution slippage, and they evaluate over too few trades to distinguish edge from
noise. The first problem is one of awareness; the second is a software-engineering
problem. \textbf{QuantEdge} addresses the second: it is a backtesting platform
that makes the rigorous path the easy one.

\subsection{Scope}

The system delivered in this submission (\texttt{v1.0.0-rc1}, tagged 2026-05-24)
provides:

\begin{itemize}
  \item Single-asset backtests across a seeded universe of 20 US equities, 5 ETFs,
        5 crypto pairs and 6 FX pairs.
  \item Daily and hourly bar resolution, each indexed against the \emph{native}
        trading calendar of its asset class (NYSE for equities and ETFs, 24$\times$5
        for FX, 24$\times$7 for crypto).
  \item Eleven strategies across five families (trend, momentum, breakout,
        mean-reversion, benchmark), each implementing a uniform signal contract.
  \item An event-driven backtest engine that enforces look-ahead-bias prevention by
        construction.
  \item Configurable commission (bps), slippage (bps or ATR-scaled), and
        position-sizing modes (fixed fraction or volatility-targeted).
  \item Standard performance dashboard: equity curve, drawdown, monthly heatmap,
        trade-P\&L scatter, rolling Sharpe, full KPI grid.
  \item Buy-and-hold and SPY benchmark overlays computed and persisted at run time.
  \item An optional, opt-in LLM-generated explanatory report grounded on the run's
        metrics.
\end{itemize}

Out of scope for this iteration: multi-asset portfolio strategies, minute-level data,
walk-forward optimisation UI, paper-trading or live-trading hooks. These are
documented in the roadmap section of \texttt{README.md} as candidates for v1.1+.

\subsection{Tech-stack rationale}

The choices reflect a preference for tools that minimise the gap between the
production-leaning quality the rubric expects and the four-week timeline the
academic calendar allows:

\begin{itemize}
  \item \textbf{Python 3.11+ / FastAPI / SQLAlchemy 2 / Pydantic v2} on the backend.
        FastAPI gives us OpenAPI for free (the live schema at \texttt{/openapi.json}
        is the source of truth for the typed frontend client). SQLAlchemy's generic
        types let the same schema run on SQLite in development and on PostgreSQL in
        production via a single \texttt{DATABASE\_URL} switch.
  \item \textbf{React 18 + TypeScript + Vite + TailwindCSS + Plotly.js} on the
        frontend. Plotly produces interactive charts without manual SVG plumbing;
        TypeScript catches schema mismatches against the OpenAPI-typed client.
  \item \textbf{pytest + Vitest} for tests. The single most important test is
        \texttt{tests/test\_engine\_no\_lookahead.py}, which injects an oracle
        strategy that ``knows'' the future close \textemdash{} the engine must still
        only let it trade at next-bar open.
  \item \textbf{LLM via a provider abstraction.} \texttt{LLMProvider} (in
        \texttt{backend/llm/base.py}) has two implementations: \texttt{NullProvider}
        (deterministic, used by tests) and \texttt{GeminiProvider} (Google Gemini,
        opt-in). LLM features are off by default.
\end{itemize}

\subsection{Four-week timeline}

The project ran from 2026-04-25 (initial commit) to 2026-05-24 (documentation pass
and submission), staffed by four contributors. Targets per week:

\begin{itemize}
  \item \textbf{Week~1 \textemdash{} 2026-04-27 to 2026-05-03.} Scaffold the
        repository, fix the directory layout, draft \texttt{requirements.txt}, agree
        on the agent boundaries.
  \item \textbf{Week~2 \textemdash{} 2026-05-04 to 2026-05-10.} First domain feature
        (TSLA as an equity asset); first LLM-generated report (proves the
        \texttt{LLMProvider} abstraction works end-to-end).
  \item \textbf{Week~3 \textemdash{} 2026-05-11 to 2026-05-17.} Frontend MVP;
        hourly bar support and per-asset-class native calendars; chart suite; PDF
        export of the AI report.
  \item \textbf{Week~4 \textemdash{} 2026-05-18 to 2026-05-24.} Data-layer
        hardening (yfinance as primary across all classes); strategy catalogue
        expansion from 5 to 11; benchmark overlays; shared crosshair and legend;
        CI workflows; full documentation pass.
\end{itemize}

The four-contributor team meant any single week's deliverable could be split across
two pairs (frontend + backend), which is reflected in the merge-PR history in
Section~\ref{sec:diary}.

\section{Project diary}
\label{sec:diary}

The four-week development cycle was punctuated by eight merged pull requests and a
handful of direct-to-\texttt{main} commits for docs and trivia. The narrative below
follows the actual order of decisions, with commit hashes given so a reviewer can
verify any claim from \texttt{git log}.

\subsection*{Week~1 \textemdash{} Foundation (2026-04-25 to 2026-04-29)}

The first commit (\texttt{4ff85c1}) bootstrapped an empty repository. Four days
later, \texttt{ecdf3bb} (``Initial project scaffold for QuantBacktest'') landed the
full directory layout we would carry through the project: a \texttt{backend/}
package with \texttt{agents/}, \texttt{api/}, \texttt{analytics/},
\texttt{backtest/}, \texttt{data/}, \texttt{database/}, \texttt{llm/},
\texttt{strategies/} subpackages; a \texttt{frontend/} folder primed for Vite; a
\texttt{tests/} folder; and \texttt{requirements.txt} pinning the backend stack.
The agent boundaries were committed to at this stage rather than discovered later
\textemdash{} we used the breakdown in the rubric (data, strategy, backtest,
analytics, orchestrator, explanation) as the starting point and held to it.

\subsection*{Week~2 \textemdash{} First domain feature + first LLM call (2026-05-08)}

Two commits on the same day represent the project's two first proofs of life. PR~\#1
(\texttt{376d30e}, ``feat: add TSLA as new equity asset'') validated the end-to-end
data path: register an asset in the seed list, run \texttt{load\_initial\_data.py},
see daily bars appear in the database. \texttt{8997d20} (``feat: add LLM generated
report'') validated the LLM seam: we wanted to make sure the \texttt{LLMProvider}
abstraction could carry a real response before we built the rest of the analytics
layer around it.

We deliberately wrote the LLM feature \emph{first} as a worst-case probe. If it
worked, every other deterministic feature would be easier; if it failed, we had time
to redesign. It worked, and the rest of the project never had to compromise its
determinism for the sake of the AI report.

\subsection*{Week~3 \textemdash{} Frontend MVP, hourly bars, charts (2026-05-12 to 2026-05-16)}

Frontend Version~1 (\texttt{1f2f5c7}) landed on 2026-05-12. The same day, PR~\#2
(\texttt{1f13139}, ``Add native calendars and hourly bar support'') reshaped the
data layer: the OHLCV composite primary key became \texttt{(asset\_id, ts,
timeframe)}, so daily and hourly bars could coexist for the same asset, and the
\texttt{OHLCVCleaner} learned to reindex onto NYSE for equities, a bespoke
24$\times$5 window for FX, and 24$\times$7 for crypto. This was the project's most
intricate single change: a missed edge case in any of the three calendars would have
silently misstated the annualisation factor for every Sharpe ratio we report.

The chart suite then went through a flurry of small fixes (\texttt{3e22d4c},
\texttt{162ceba}, \texttt{ab1e2de}) that brought the strategy-category badges,
chart-tooltip alignment, and sidebar navigation up to standard. \texttt{a3b6635}
(``Chunk upserts, add LLM retries, and fix time indexes'') addressed the SQLite
parameter-count ceiling we hit when bulk-loading multi-year hourly history.

Late in the week, the PDF-export feature shipped (\texttt{0a7edb1}). The choice to
render the report client-side via \texttt{jsPDF} rather than server-side via a
headless browser kept the backend dependency footprint small and avoided the need
for the OS-level Chromium libraries that other PDF approaches require.

\subsection*{Week~4 \textemdash{} CI hardening, data layer, strategies, benchmarks (2026-05-19 to 2026-05-22)}

The week opened with two Claude Code GitHub Actions workflows: \texttt{claude.yml}
(\texttt{3d50b41}) for \texttt{@claude} mentions on issues and PRs, and
\texttt{claude-code-review.yml} (\texttt{6a38f60}) for an automatic
\texttt{/code-review} pass on every PR. The team's view of generative AI as a
collaborator-not-author shaped both workflows: they only respond to explicit human
triggers and they never self-merge.

Build-artefact hygiene then dominated two days (\texttt{4ce0311},
\texttt{820d610}). TypeScript compile outputs had begun leaking into
\texttt{frontend/src/}, which would have polluted every diff. PR~\#4 added the
\texttt{check-no-build-artifacts.yml} workflow that fails the build if any
\texttt{.js}, \texttt{.d.ts}, or stale \texttt{.tsbuildinfo} appears where it
should not.

The data layer consolidated on yfinance as the primary source across every asset
class (PR~\#5, \texttt{c2b6cce}). The Binance \texttt{ccxt} fallback remained for
crypto hourly windows older than the 730-day yfinance cap, and Stooq remained for
FX EOD fallback. Fixed-fractional sizing was corrected for crypto and FX (PR~\#6,
\texttt{ae082ef}).

PR~\#7 (\texttt{b0676ed}, ``added 6 new strategies'') doubled the strategy
catalogue to eleven by adding MACD Crossover, Ichimoku Cloud, Keltner Channels,
Time Series Momentum, Stochastic Oscillator, and CCI. Each was added to the
registry, exposed via JSON Schema so the frontend form rendered automatically, and
covered by a unit test that asserts the signal series on a known fixture.

PR~\#8 (\texttt{c0801c7}, ``Support benchmark overlays (API + UI)'') and its
sibling commits (\texttt{be6cf8a}, \texttt{55fa4b0}, \texttt{d432105}) brought the
chart suite to its final polish: a same-asset buy-and-hold and an SPY buy-and-hold
curve are now computed and persisted at run time, every chart shares a single
legend that toggles them, and a shared crosshair tooltip keeps values aligned
when the mouse hovers anywhere on the timeline.

The week closed with \texttt{d1efa3f} (Quick-Start docs) and \texttt{4a9e6ae}
(yfinance/httpx/curl\_cffi version bumps). The yfinance bump from 0.2.46 to 1.3.0
required adding \texttt{curl\_cffi} as a transitive dependency \textemdash{} a
small change that we noted in \texttt{CHANGELOG.md} because it could surprise
contributors who ran into install issues on locked-down corporate machines.

\subsection*{Documentation pass (2026-05-24)}

\texttt{79e200a} (``Updated all the current documentation files'') refreshed the
existing docs \textemdash{} \texttt{README.md}, \texttt{ARCHITECTURE.md},
\texttt{CLAUDE.md}, \texttt{ONBOARDING.md}, and the four
\texttt{docs/\{agents,calendars,data-sources,strategies\}.md} \textemdash{} against
the final code state. A second pass on the same day created the artefacts the
rubric requires (this PDF among them): \texttt{AGENTS.md}, \texttt{CHANGELOG.md},
\texttt{CONTRIBUTING.md}, \texttt{CITATIONS.md}, \texttt{AI\_USAGE.md},
\texttt{SECURITY.md}, \texttt{CODE\_OF\_CONDUCT.md},
\texttt{docs/user-guide.md}, \texttt{docs/api.md},
\texttt{docs/architecture-diagrams/}, and this \texttt{docs/academic/} bundle.

\subsection*{Post-submission CI automation (2026-05-24)}

On the same evening as the submission, PR~\#9 (\texttt{612e39e}, ``chore(ci): add
nightly headless docs-refresh workflow'') introduced \texttt{nightly-docs.yml}: a
GitHub Actions job that wakes at 02:00~UTC each night, runs Claude Code in headless
mode with a read-only Bash allowlist, and opens a pull request whenever it finds
documentation drift against the last 24~hours of commits. The workflow never touches
\texttt{main} directly.

Two small fixes followed immediately. \texttt{4dfc948} switched the authentication
method to \texttt{CLAUDE\_CODE\_OAUTH\_TOKEN} (the token type the
\texttt{claude-code-action} expects), and PR~\#10 (\texttt{58a4811}) granted the
\texttt{id-token: write} permission that the action requires to mint a GitHub OIDC
token at startup.

\subsection*{Node 24 Actions upgrade (2026-05-24)}

GitHub announced the retirement of Node 20--based runner actions on 2026-06-02.
Commit \texttt{01d9c30} (PR~\#12) bumped the three flagged actions
(\texttt{actions/checkout}, \texttt{actions/setup-python}, \texttt{actions/setup-node})
to their v6 majors across all four workflows in a single pass. The permissions comment
in \texttt{nightly-docs.yml} was updated to record why \texttt{id-token: write} is
required: \texttt{claude-code-action} mints a GitHub OIDC token at startup regardless
of which model-auth path is active, so the permission is not optional.

\subsection*{Things we tried and abandoned}

A few decisions did not survive the four weeks:

\begin{itemize}
  \item \textbf{Asynchronous backtest submission.} The first design queued runs and
        polled status via WebSocket. We backed it out because synchronous \texttt{POST}
        worked well enough at v1 scale and removed an entire layer of state to debug.
        The Dashboard still polls \texttt{GET /backtests} every five seconds so it
        will surface any future async-submission feature without UI changes.
  \item \textbf{Postgres in development from day one.} We initially set up Docker
        Compose with Postgres but it slowed local iteration noticeably. SQLite became
        the default and the schema was kept on SQLAlchemy generic types so a
        \texttt{DATABASE\_URL} swap reaches Postgres without code changes.
  \item \textbf{LLM tool-use via the provider SDK's native function calling.} The
        Orchestrator Agent now parses tool calls from a JSON envelope in the model
        output rather than relying on Gemini's native function-calling API, because
        the latter is not directly portable to other providers. The cost is a
        small parsing step; the benefit is provider-agnostic tool dispatch.
\end{itemize}

\section{Methodology: economics, mathematics, data science}
\label{sec:methodology}

This section makes the project's analytical contract precise: the event-driven
backtest semantics, the KPI formulas with their per-asset-class annualisation
factor, the strategy taxonomy, and the LLM-provider abstraction.

\subsection{Event-driven backtest semantics and look-ahead bias}
\label{sec:lookahead}

Look-ahead bias is the canonical pitfall of backtest construction: a strategy
``sees'' information that, in real time, it would not have had access to when it
needed to act \cite{lopezdeprado2018,bailey2014}. The most common form is a
\emph{single-bar} leak \textemdash{} a strategy reads bar~$t$'s closing price and
then places an order that fills at bar~$t$'s open or close. Across a long backtest
the false edge accumulates and inflates every downstream metric. QuantEdge
forbids the leak architecturally: orders generated during bar~$t$ fill at the
\emph{open of bar~$t{+}1$}, with commission and slippage applied to the fill
price. The rule is implemented once, in \texttt{backend/backtest/engine.py}, and
asserted by an oracle-strategy test in \texttt{tests/test\_engine\_no\_lookahead.py}
that injects a strategy that ``knows'' the future close; the test passes only if
the engine still forces the oracle to trade at next-bar open.

In pseudocode the engine loop is:

\begin{lstlisting}[language=Python]
for t in range(len(bars)):
    # Target position from data up to and including bar t.
    target = strategy.generate_signals(bars[:t+1]).iloc[t]   # in {-1, 0, 1}
    if t == len(bars) - 1: break       # no t+1 open to fill at
    delta = target - portfolio.current_position
    if delta == 0: continue
    # Fill at bar (t+1) open, applying slippage and commission.
    fill_price = bars.iloc[t+1].open * (1 + sign(delta) * slippage_bps / 10_000)
    commission = abs(delta * portfolio.size) * fill_price * commission_bps / 10_000
    portfolio.apply(delta, fill_price, commission)
\end{lstlisting}

The signal contract is uniformly \texttt{pd.Series[int]} with values in
$\{-1, 0, +1\}$ for ``short, flat, long''. Strategies do not multiply their own
signals by a position-size factor; the engine handles sizing via either
\emph{fixed fractional} (a fraction of equity per trade) or \emph{volatility
targeting} (position scaled inversely to realised volatility). An optional
circuit-breaker forces the run flat if equity draws down beyond a user-supplied
threshold.

\subsection{Performance metrics}
\label{sec:kpi}

The KPI definitions follow the conventions in Bacon~\cite{bacon2008}. Let
$\{r_t\}_{t=1}^{T}$ denote the per-bar returns of the equity curve
$\{V_t\}_{t=0}^{T}$, $k$ the number of bars in one calendar year
(Section~\ref{sec:calendars}), and $y = T/k$ the run's length in years.

\paragraph{Total return and CAGR.}
\begin{equation}
  R = \frac{V_T}{V_0} - 1,\qquad
  \mathrm{CAGR} = \left(\frac{V_T}{V_0}\right)^{1/y} - 1.
\end{equation}

\paragraph{Annualised volatility.}
\begin{equation}
  \sigma_{\text{ann}} = \sqrt{k}\,\sigma(r_t).
\end{equation}

\paragraph{Sharpe and Sortino ratios.}
\begin{equation}
  \mathrm{Sharpe} = \sqrt{k}\,\frac{\bar r}{\sigma(r_t)},\qquad
  \mathrm{Sortino} = \sqrt{k}\,\frac{\bar r}{\sigma_{-}(r_t)},
\end{equation}
where $\sigma_{-}(r_t)$ is the downside semi-deviation $\sqrt{\mathbb{E}[\min(r_t, 0)^2]}$.
We adopt the zero-risk-free-rate convention for retail backtesting
\cite{sharpe1966,sortino1994}.

\paragraph{Maximum drawdown and Calmar ratio.}
\begin{equation}
  \mathrm{maxDD} = \min_{t} \frac{V_t - \max_{s \le t} V_s}{\max_{s \le t} V_s},\qquad
  \mathrm{Calmar} = \frac{\mathrm{CAGR}}{|\mathrm{maxDD}|}.
\end{equation}
Calmar~\cite{young1991} captures the pain-adjusted return: how much annual growth
the strategy delivered per unit of drawdown the user had to endure.

\paragraph{Trade-level metrics.} Let $\{p_i\}_{i=1}^{N}$ denote per-leg net P\&L:
\begin{equation}
  \mathrm{winrate} = \frac{|\{i:p_i > 0\}|}{N},\qquad
  \mathrm{PF} = \frac{\sum_{p_i > 0} p_i}{\sum_{p_i < 0} |p_i|},
\end{equation}
\begin{equation}
  \mathrm{expectancy} = \mathrm{winrate}\cdot \overline{p_{+}} - (1-\mathrm{winrate}) \cdot |\overline{p_{-}}|,
\end{equation}
where $\overline{p_{+}}$ and $\overline{p_{-}}$ are the means of the positive and
negative P\&L populations respectively~\cite{pardo2008,tharp1998}.

\subsection{Annualisation factor by asset class and timeframe}
\label{sec:calendars}

A Sharpe ratio computed against the wrong $k$ misstates the annualised number by
tens of percent. QuantEdge looks up $k$ from
\texttt{backend/analytics/periods.py} as a function of \texttt{(asset\_class,
timeframe)}:

\begin{table}[h]
\centering
\small
\begin{tabular}{lcc}
\toprule
\textbf{Asset class} & \textbf{Daily $k$} & \textbf{Hourly $k$} \\
\midrule
Equity / ETF (NYSE) & 252  & 1\,638 \\
FX (24$\times$5)    & 260  & 6\,240 \\
Crypto (24$\times$7)& 365  & 8\,760 \\
\bottomrule
\end{tabular}
\caption{Annualisation factor $k$ per asset class and timeframe. NYSE hourly uses
$252 \times 6.5$ elapsed market hours per session.}
\end{table}

The NYSE hourly value reflects elapsed market hours, not the literal bar count,
because the first hourly bar of a regular session covers only the opening
30-minute auction. Early-close sessions (about ten per year) further reduce the
count slightly; the lookup is calibrated against the
\texttt{exchange\_calendars} XNYS schedule rather than computed analytically.

\subsection{Strategy taxonomy}

The eleven shipped strategies populate five families. Each family responds to a
different market regime; running them side-by-side makes that regime sensitivity
visible, which is the explicit pedagogical motive for shipping the breadth.

\begin{itemize}
  \item \textbf{Trend} (SMA Crossover \cite{pardo2008}, MACD Crossover \cite{appel2005},
        Ichimoku Cloud \cite{hosoda1969}): designed to capture sustained directional moves.
  \item \textbf{Momentum} (Time Series Momentum \cite{moskowitz2012}): own-asset trailing-return
        sign; profits from 1\textendash{}12 month persistence.
  \item \textbf{Breakout} (Donchian \cite{faith2007}, Keltner \cite{keltner1960}):
        long on channel breaks; trades clean trending regimes.
  \item \textbf{Mean reversion} (Bollinger \cite{bollinger2001}, CCI \cite{lambert1980},
        RSI \cite{wilder1978}, Stochastic Oscillator \cite{lane1984}): trade fades within
        range-bound regimes; bleed in trends.
  \item \textbf{Benchmark} (Buy and Hold): always long; serves as the reference both
        for active strategies and for the on-chart same-asset overlay.
\end{itemize}

All eleven strategies share the same signal contract and therefore the same engine
path. Adding a twelfth is a single-file change plus a registry edit; the frontend's
parameter form regenerates from the new JSON Schema.

\subsection{The LLM-provider abstraction}

Two of the six runtime agents \textemdash{} the Orchestrator and the Explanation
Agent \textemdash{} are LLM-backed, and every call they make goes through
\texttt{LLMProvider.generate(\ldots)}. Two implementations ship:
\texttt{NullProvider} (deterministic canned text) and \texttt{GeminiProvider}
(Google's Gemini API). The abstraction earns its keep three ways: tests use
\texttt{NullProvider} exclusively and stay reproducible without an API key;
switching provider is a config change, not a code change; and it makes the
project's policy explicit \textemdash{} \emph{the deterministic surface is the
default}, and any non-determinism is a conscious operator opt-in
(\texttt{LLM\_ENABLED=true} plus a populated \texttt{GEMINI\_API\_KEY}).

\section{Sample results and economic reflection}
\label{sec:results}

This section reports a representative backtest, contextualises it against the two
benchmark overlays the system computes automatically, and reflects on what the
result implies about the strategy's regime sensitivity. The numeric cells in
Table~\ref{tab:results} are populated by the team from a run they execute
themselves, so the table represents the actual code path rather than a fabricated
example.

\subsection{Run methodology}

\paragraph{Configuration.} We report a backtest of the \texttt{sma-crossover}
strategy on \texttt{SPY} (SPDR S\&P~500 ETF Trust), over 2020-01-01 to 2024-12-31,
daily bars (\texttt{timeframe="1d"}). Strategy parameters:
\texttt{fast\_window=20}, \texttt{slow\_window=50}, \texttt{allow\_short=false}.
Execution parameters: \texttt{initial\_cash=10\,000}, \texttt{commission\_bps=5},
\texttt{slippage\_bps=2}, \texttt{sizing\_mode="fixed\_fraction"},
\texttt{risk\_fraction=1.0}, no max-DD circuit breaker. The annualisation factor
is $k=252$ (NYSE daily, Section~\ref{sec:calendars}).

\paragraph{Reproducibility.} The exact API call:

\begin{lstlisting}[language=bash]
curl -s -X POST http://127.0.0.1:8000/backtests \
  -H 'Content-Type: application/json' \
  -d '{
    "asset_symbol": "SPY",
    "strategy_slug": "sma-crossover",
    "start_date":   "2020-01-01T00:00:00",
    "end_date":     "2024-12-31T00:00:00",
    "params":       {"fast_window": 20, "slow_window": 50, "allow_short": false},
    "timeframe":    "1d"
  }'
\end{lstlisting}

Then \texttt{GET /backtests/\{run\_id\}/metrics} and the two
\texttt{/benchmark/\{kind\}/metrics} endpoints supply the table values.

\subsection{Results table}

\todo[inline]{TEAM: replace the \emph{TBD} cells in Table~\ref{tab:results} below
with the numbers your actual run returned. The table compiles either way; the
TBDs make the gap visible to a reviewer.}

\begin{table}[ht]
\centering
\small
\renewcommand{\arraystretch}{1.2}
\begin{tabular}{lccc}
\toprule
\textbf{Metric} & \textbf{SMA Crossover} & \textbf{Same-asset B\&H} & \textbf{SPY B\&H} \\
\midrule
Total return       & \textit{TBD} & \textit{TBD} & \textit{TBD} \\
CAGR               & \textit{TBD} & \textit{TBD} & \textit{TBD} \\
Annualised vol     & \textit{TBD} & \textit{TBD} & \textit{TBD} \\
Sharpe             & \textit{TBD} & \textit{TBD} & \textit{TBD} \\
Sortino            & \textit{TBD} & \textit{TBD} & \textit{TBD} \\
Max drawdown       & \textit{TBD} & \textit{TBD} & \textit{TBD} \\
Calmar             & \textit{TBD} & \textit{TBD} & \textit{TBD} \\
Trade count        & \textit{TBD} & 1 & 1 \\
Win rate           & \textit{TBD} & --- & --- \\
Profit factor      & \textit{TBD} & --- & --- \\
\bottomrule
\end{tabular}
\caption{KPIs for the representative SMA-crossover run, alongside the two
benchmark overlays computed automatically by \texttt{BacktestAgent}. Trade-block
cells are dashed for benchmarks (buy-and-hold has only one entry and one exit).
Cells marked \emph{TBD} are filled in by the team from the actual run.}
\label{tab:results}
\end{table}

\subsection{Economic reflection}

The economic story we expect Table~\ref{tab:results} to tell follows from the
nature of the strategy and the chosen window. SMA Crossover is a trend-follower:
it goes long when the 20-day SMA crosses above the 50-day SMA and exits when it
crosses back. The 2020\textendash{}2024 window spans the COVID drawdown (March
2020), the equity bull market of 2020\textendash{}2021, the 2022 bear market, and
the 2023\textendash{}2024 recovery \textemdash{} four distinct regimes inside five
years. A 20/50 SMA crossover should capture the trends of 2021 and 2023\textendash{}2024
cleanly, sit in cash during much of 2022, but also lose money to whipsaws around
the regime transitions (early 2020 and mid-2022).

We therefore expect three observations to hold once the team populates the table:

\begin{enumerate}
  \item \textbf{Lower CAGR than SPY buy-and-hold} but \textbf{noticeably lower
        maximum drawdown}, because the strategy goes flat during 2022's bear
        market. The pain-adjusted comparison (Calmar) is the more interesting
        head-to-head.
  \item \textbf{A Sharpe ratio close to buy-and-hold's} \textemdash{} the
        strategy's higher hit rate during clear trends is offset by the
        cost of being in cash through the recoveries.
  \item \textbf{A trade count modest enough to demand caution.} Five years of daily
        bars typically produce 10\textendash{}30 trades for a 20/50 SMA system, which
        is below the rule-of-thumb threshold of $\sim$100 trades for statistical
        significance. Any strong single-asset result here should be re-tested on
        several assets before being trusted.
\end{enumerate}

\todo[inline]{TEAM: replace the three-point list above with the actual observations
once Table~\ref{tab:results} is populated. The numbered list is a template; rewrite
it in your own words to reflect the run you actually executed.}

\subsection{Limitations of the demonstration}

The result above is a single asset over a single window with a single parameter
set. It is illustrative, not predictive. Three caveats apply:

\begin{itemize}
  \item \textbf{Survivorship bias} \textemdash{} yfinance carries only currently
        listed tickers, so any equity universe is biased upward.
  \item \textbf{Single-window evaluation} \textemdash{} the rubric's preference for
        walk-forward / out-of-sample evaluation is correct; SMA-crossover's
        in-sample fit on a single window does not establish edge.
  \item \textbf{The README's retail-trader tip applies} \textemdash{} if the
        backtested maximum drawdown is 20\%, the user should mentally prepare for
        30\% in live trading, both because of regime change and because realistic
        slippage on retail accounts is usually worse than the 2-bps default.
\end{itemize}

\section{Lessons learned}
\label{sec:lessons}

\subsection{What worked}

\paragraph{Pydantic-driven JSON Schemas as the contract between backend and frontend.}
Every strategy in \texttt{backend/strategies/} declares a Pydantic \texttt{Config}
class; \texttt{GET /strategies} serialises it as a JSON Schema; the React form
component renders a complete parameter UI from that schema with no per-strategy
frontend code. Adding the sixth-through-eleventh strategies (PR~\#7) cost zero
frontend work as a result \textemdash{} the single biggest productivity
multiplier in the project.

\paragraph{The six-agent abstraction.} Splitting the runtime into named agents,
each inheriting \texttt{BaseAgent[TIn, TOut]} with a single public
\texttt{run(payload)} method, gave us seam-by-seam testability and a natural
docking point for the LLM orchestrator. The boundary between the two LLM-backed
agents and the four deterministic agents was the right boundary: it kept the test
suite reproducible while leaving the LLM path open as an opt-in feature.

\paragraph{The $\{-1, 0, +1\}$ signal contract.} Forbidding strategies to manage
their own position size, commissions, or fills made the engine the single point of
truth for the bar-$t \to$ bar-$t{+}1$ rule, and meant that the
look-ahead-bias oracle test in \texttt{tests/test\_engine\_no\_lookahead.py} is the
only invariant guard we need. Every strategy passes through the same fill code
path; a bug anywhere in fill semantics surfaces everywhere at once.

\paragraph{Native trading calendars per asset class.} The temptation to reindex
every asset onto NYSE for ``consistency'' would have silently misstated the
annualisation factor for FX and crypto by 5\textendash{}50\%. Building
calendar-aware reindexing into the cleaner from the start (PR~\#2) cost a few
hours and saved us from a class of subtle correctness bugs.

\subsection{What did not work}

\paragraph{Yahoo Finance's 730-day hourly cap.} The cap is a hard external
constraint that propagates into the UI (date pickers had to gain a per-timeframe
maximum-look-back constraint) and into the data layer (the Crypto Fetcher learned
to fall back to Binance via \texttt{ccxt}). We absorbed the cap; the alternatives
(paid intraday data, scraping) cost more than the limitation imposes.

\paragraph{LLM tool-call JSON parsing.} The Orchestrator Agent parses tool calls
out of the model's response by extracting a JSON envelope. This is fragile:
models occasionally wrap the JSON in markdown fences or prose. We chose
provider-portability over robustness for v1; a future iteration could fall back
to provider-native function calling where available.

\paragraph{Asynchronous backtest execution.} The first design had submission
return immediately and the UI poll for status; we backed it out to a synchronous
\texttt{POST /backtests} because runs on daily bars complete in under a second
and the async machinery added an entire state-management layer that nothing in
v1 actually needed (Section~\ref{sec:diary}).

\subsection{What we would do differently}

\begin{itemize}
  \item \textbf{Walk-forward UI from day one.} The Strategy Agent already exposes a
        chronological train/test split, but the UI does not surface it. Without a UI
        affordance, users will tune parameters by re-running the full backtest
        \textemdash{} the canonical recipe for overfitting~\cite{bailey2014}. The
        UI for walk-forward should have been a v1 feature, not a v1.1 candidate.
  \item \textbf{Postgres as the default development database.} SQLite was fast to set
        up but does not catch the subset of bugs that only show up under concurrent
        writes or under strict referential integrity. The schema is already
        Postgres-compatible (SQLAlchemy generics only); next time we would default
        to Postgres and let users opt down to SQLite.
  \item \textbf{Earlier CI hygiene.} The \texttt{check-no-build-artifacts.yml}
        workflow should have been the first \texttt{.github/workflow} we wrote, not
        the third. Build-artefact leakage produced noisy diffs for a week before we
        formalised the guard.
  \item \textbf{Smaller, more frequent PRs.} The 32 development-cycle commits landed
        via only 8 merged PRs; some touched a dozen files across two subsystems.
        Reviewers (human or AI) work better on small, focused diffs.
\end{itemize}

\subsection{What carries forward}

The patterns we would copy verbatim into the next project: the
schema-as-API discipline, the agent abstraction with a single public
\texttt{run} method, the deterministic-by-default LLM stance with
\texttt{NullProvider} for tests, the per-asset-class trading-calendar layer, the
$\{-1, 0, +1\}$ signal contract, and the ``\emph{Last verified against code}''
footer convention on every documentation file.

\section{Acknowledgement of generative-AI tools}
\label{sec:ai}

This section is required by the course rubric and is a compressed version of
\texttt{AI\_USAGE.md} at the repository root, which carries the full per-tool table
and policy statement.

\subsection*{Team policy}

The team used generative-AI tools as collaborators, not as authors. Three
durable rules governed every contribution:

\begin{enumerate}
  \item Every line of code that lands in \texttt{main} is reviewed by a human. AI
        output that no team member can explain is not merged.
  \item No AI tool is given access to credentials, API keys, or production data.
        The \texttt{.env} file is gitignored; only \texttt{.env.example} (no secrets)
        ships in the repository.
  \item The LLM features inside the product itself are opt-in:
        \texttt{backend/config.py} defaults to \texttt{LLM\_ENABLED=false}, the
        test suite uses the deterministic \texttt{NullProvider}, and activating
        the live \texttt{GeminiProvider} requires three environment variables.
\end{enumerate}

We take collective responsibility for the whole repository, including text and
code that was originally drafted by an AI tool and then reviewed and accepted.

\subsection*{Tools used}

\textbf{Claude Code} (Anthropic, CLI agent) was the primary coding assistant for
multi-file refactors, agent scaffolding, documentation drafts, and the bulk of the
LaTeX source for this submission; it left footprints in the repository via the
\texttt{.github/workflows/claude.yml} (issue/PR \texttt{@claude} mention handler)
and \texttt{.github/workflows/claude-code-review.yml} (per-PR \texttt{/code-review})
workflows. \textbf{Claude.ai} (Anthropic, web) was used by team members for design
discussions and prose drafting. \textbf{ChatGPT} (OpenAI) was used opportunistically
for isolated syntax checks and LaTeX-table formatting. \textbf{GitHub Copilot}
(GitHub/OpenAI) provided in-editor autocomplete for routine code; one merge-resolution
commit (\texttt{b9fc033}) was authored by the \texttt{copilot-swe-agent[bot]} account
as a one-off, not a sustained authoring pattern. \textbf{Cursor} (Anysphere) was used
by one team member as their primary editor for some of the frontend work.
\textbf{Google Gemini} (web) was used to spot-check factual claims about NYSE hours
and \texttt{ccxt} pagination semantics; \textbf{Google Gemini API} is also wired into
the product itself as the optional live \texttt{LLMProvider}, but that is a product
feature, not a development tool.

The team did \emph{not} use Codex, Replit Ghostwriter, Sourcegraph Cody, Tabnine,
Codeium, or Amazon CodeWhisperer.

\subsection*{Where AI was not used}

The strategy formulas in \texttt{backend/strategies/*.py} were transcribed from the
academic and practitioner sources cited in \texttt{CITATIONS.md} and
Section~\ref{sec:methodology}; the team independently re-derived each one before
implementation. The look-ahead-bias enforcement in
\texttt{backend/backtest/engine.py} and its oracle-strategy test were designed and
written by the team. The database schema, the API contract, and the four-week
project timeline were team decisions, not AI-generated.

\subsection*{Reproducibility caveat}

AI sessions are not natively reproducible across runs or prompt phrasings. We
mitigate this via diff-level human review, the test suite's coverage of the
critical correctness invariants, and prompt summaries in the project diary
(Section~\ref{sec:diary}) for the most impactful changes.


% ---------- Bibliography ----------------------------------------------------
\printbibliography[heading=bibnumbered,title={References}]

\end{document}
